\RequirePackage{etex}

% Vi benytter 2 sidet, a4 størrelse og dansk sprog. 
\documentclass[11pt,a4paper,oneside,danish]{memoir}
\usepackage[paper=A4,pagesize]{typearea}
 \usepackage[utf8]{inputenc} % danske tegn

% Choose language by uncommenting danish or english
\usepackage[danish]{babel} % dansk sprogpakke
% \usepackage[english]{babel}


%% Pakker til fonts, farver, referencer og ligenende er placeret relevant sted i dokumentet.
     
     %% FONT
     \usepackage[default,osfigures,scale=0.95]{opensans} % Alternatively use the option 'defaultsans' instead of 'default' to replace the sans serif font only.
     \usepackage[T1]{fontenc}
     
     % Common packages 
     \usepackage{footnote} % Package used to controle footnote(specially in tables etc.)
     \usepackage{amsmath}
     \usepackage{slantsc}
     \usepackage{multirow} % Merge table rows
     \usepackage{longtable} % Table that can span multiple pages
     \usepackage{geometry}
     \usepackage[section]{placeins} % \FloatBarrier
     \usepackage[compatibility=false]{caption}
     \captionsetup{font=small,labelfont=bf,textfont=it}
     \usepackage{graphicx}
     \usepackage{afterpage}
     \usepackage[draft]{fixme} %hvor 'options' bl.a. inkluderer final, draft, danish, inline og margin
     \usepackage{rotating}
     \usepackage{eso-pic} % til inds�ttelse af forside.
     \usepackage{rotating} % rotate stuff(EVERYTHING!11)

	% % MATH STUFF 
    %\usepackage{algorithmic}
    %\usepackage{harpoon} %Tilf�jer harpun-pile over st�rrelser i math, m.m.
    %\usepackage{pst-plot} % For axes   
    %\usepackage{breqn}
    %\usepackage{fourier}  %................... Roman+math - Utopia
    
  
    %\usepackage{nomentbl} % http://ctan.mackichan.com/macros/latex/contrib/nomentbl/nomentbl.pdf

    %\usepackage{psfrag} % insæt tag, mest brugt i forhold til grafer

%%% DIVERSE CUSTOMS PAKKET NED I FILER FOR OVERBLIK
   % \input{styles/HaxTilPreamble.tex}
   % \input{styles/CustomFunctions.tex}
%	\input{Styles/code.tex} 
    \usepackage{todonotes}


% new commands for individually colored todo 
\newcommand{\todosander}[2]{
\todo[#1, color=green!40]{#2 \\- Sander}
}
\newcommand{\todosimon}[2]{
\todo[#1, color=red!40]{#2 \\- Simon}
}
\newcommand{\todorasmus}[2]{
\todo[#1, color=yellow!40]{#2 \\- Rasmus}
}
\newcommand{\todokristoffer}[2]{
\todo[#1, color=purple!40]{#2 \\- Simon}
}


% Color package. Use with listings package and use for color a table
%\usepackage[usenames,dvipsnames,svgnames,table]{xcolor}

    %%%%%%%%%%%%%%%%%%%%%%%%%%%%%%%%%%%%%%%%%%%%%%%%%%%%%%%%%%%%%%%%%%%%%%%%%%%%%
%                            Custom for funktioner                          %
%                     			      	        		                    %
%%%%%%%%%%%%%%%%%%%%%%%%%%%%%%%%%%%%%%%%%%%%%%%%%%%%%%%%%%%%%%%%%%%%%%%%%%%%%


% Package for wrapping figures in text
\usepackage{wrapfig}

% Set the path where LaTeX looks for pictures.
\graphicspath{{figure/}}

% You need a newsubfloat element to use subcaption
\newsubfloat{figure}

% Command to set caption styles
\captionnamefont{\bfseries\small}
\captiontitlefont{\itshape\small}
\subcaptionlabelfont{\bfseries\small}
\subcaptionfont{\itshape\small}

% When we don't use chapter numbering, 
% the figure numbering are shown like 
% fx 0.3 if we don't use this line.
\counterwithout{figure}{chapter}


% Command \figur{filename}{caption}{label}{width} 
% for inserting a new figure.
\newcommand{\figur}[4]{ 
 \begin{figure}[ht] 
 \centering 
 \includegraphics[width=#4\textwidth]{#1} 
 \caption{#2} 
 \label{#3} 
 \end{figure}
 } 
 
% Command \wrapfigur{filename}{caption}{label}{width}{r/l}
% for inserting a new figure wraped in the text. 
\newcommand{\wrapfigur}[5]{ 
 \begin{wrapfigure}{#5}{#4 \textwidth}
   \begin{center}
    \includegraphics[width=#4\textwidth]{#1}
   \end{center}
   \caption{#2}
   \label{#3} 
 \end{wrapfigure}
}
 
% Command \subfigur{filename1}{filename2}{caption}
% {subcaption1}{subcaption2}{label1}{label2}
% {width1}{width1} for inserting a new figure.
\newcommand{\subfigur}[9]{
\begin{figure}[ht]
\centering
\subbottom[{#4}\label{#6}]%
    {\includegraphics[width=#8\textwidth]{#1}}\hfill
\subbottom[{#5}\label{#7}]%
    {\includegraphics[width=#9\textwidth]{#2}}
\caption{#3}
\label{fig:subfigures}
\end{figure}
}
% Command \truefigure{filename}{caption}{label}
% for inserting a picture without scaling.
\newcommand{\truefigure}[3]{ 
 \begin{figure}[ht] 
 \centering 
 \includegraphics{#1} 
 \caption{#2} 
 \label{#3} 
 \end{figure}
 } 

% Keeps floats in the section. 
% Help control place to put pix
\usepackage[section]{placeins}

% Commands for rightaligning wraped text in a tabular
\newcommand{\rr}{\raggedright}
\newcommand{\tn}{\tabularnewline}

\usepackage{multirow}

%Tabular pakke som Wex har fundet, men som vi ikke bruger alligvel
%\usepackage{tabu}

% Include code pieces from files
% Documentation:
% ftp://ftp.tex.ac.uk/tex-archive/macros/latex/contrib/listings/listings.pdf
\usepackage{listings}


%Package with enumirate command with extra options (used within cells of tabular)
\usepackage[inline]{enumitem}
    
    % add a box using the themecolor. Can be used to wrap around a chapintro for introducing an assignment
    %use \begin{assignment} and end\{assignment} to box it
    %http://texdoc.net/texmf-dist/doc/latex/mdframed/mdframed.pdf
     
    \usepackage[xcolor]{mdframed} 
    \newmdenv[linecolor=ThemeColor]{assignment}
%    \usepackage{Placeins}
    
% % % % % % % % % % % % % % ********% % %Generelt design % % % % % % % % % % % % % % % % % % %
 
%\usepackage[table,dvipsnames]{xcolor} % 'table' er i stedet for \usepackage{colortbl} og 'dvipsnames' er et farveskema hvor man angiver farver med navne

\usepackage{xcolor} % farver kan indsætttes i henhold til RGB med definecolor

%Valg af temafarve % Denne farve bruges i chapterstyle %
\definecolor{ThemeColor}{RGB}{209, 67, 26}

% Valg af font %
\usepackage[sc]{mathpazo}
\linespread{1.05}         % Palatino needs more leading (space between lines)
\usepackage[T1]{fontenc}

% Pakker inkludere overordnede designvalg. Pakkerne er valgt for overskuelighed %
	 %%%%%%%%%%%%%%%%%%%%%%%%%%%%%%%%%%%%%%%%%%%%%%%%%%%%%%%%%%%%%%%%%%%%%%%%%%%%%
%                             Chapter and Section                           %
%                     			      Style	       		                    %
%%%%%%%%%%%%%%%%%%%%%%%%%%%%%%%%%%%%%%%%%%%%%%%%%%%%%%%%%%%%%%%%%%%%%%%%%%%%%


% % % % Redefined so that header will show even on chapter start. % % % %
\makeatletter
    \let\stdchapter\chapter
    \renewcommand*\chapter{%
    \@ifstar{\starchapter}{\@dblarg\nostarchapter}}
    \newcommand*\starchapter[1]{%
        \stdchapter*{#1}
        \thispagestyle{fancy}
        \markboth{\MakeUppercase{#1}}{}
    }
    \def\nostarchapter[#1]#2{%
        \stdchapter[{#1}]{#2}
        \thispagestyle{fancy}
    }
\makeatother




\usepackage{titlesec} % Pakke til opsætning af chapter og sektioner


% % defining new font for the chapternumber. Calling it FuckingHuge 
\newcommand{\FuckingHuge}{% Needs a 'stretchable' font
      \usefont{\encodingdefault}{\rmdefault}{b}{n}%
      \fontsize{35}{45}\normalfont\color{ThemeColor}} %60 80


%%% TAL I OVERSKRIFTER, både chapters og sections. 
    \maxsecnumdepth{subsubsection} %subsubsection er dybdeniveau for numering.
   	\setsecnumdepth{subsubsection}


  
 % % % % CHAPTER OPSÆTNING % % % % % % % % % % % % % % % % % % % % %

%Hang thoose numbers! I margin med tallene for både sections og chapters! En lækker chapter style!
\chapterstyle{hangnum}
 
%Modificering af ovenstående style
%Placering af chapter overskriften på siden i forhold til sidens start %
\renewcommand*{\chapterheadstart}{\vspace*{-20pt}}


% FARVE PÅ CHAPTER % %
%TItel farve og font valg
\renewcommand*{\chaptitlefont}{\normalfont\Huge\color{ThemeColor}}
%Numering farve og font valg
\renewcommand{\chapnumfont}{\normalfont\FuckingHuge\color{ThemeColor}}



% % % % SECTION OPSÆTNING % % % % % % % % % % % % % % % % % % % %
\hangsecnum % vi smidder numering ud i marken

%Vi benytter tema farven til numeringen
\renewcommand\thesection{\color{ThemeColor}\thechapter.\arabic{section}}



    
    


	 \input{Styles/header.tex}
	 
% Setup the margins:
\setlrmarginsandblock{1in}{*}{1.5}
\setulmarginsandblock{3cm}{4cm}{*}

% Setup header and footer: %
% \setheadfoot{headheight}{footskip}
% \setheaderspaces{headdrop}{headsep}{ratio}
\setheadfoot{45pt}{1cm}
\setheaderspaces{1cm}{*}{*}
\checkandfixthelayout
% Se dette pdf for dokmentaton:
% www.tug.dk/_media/foredrag/2003-12-03/slides-src.pdf

 \checkandfixthelayout[nearest]
    \checkthelayout
    \fixthelayout
    \baselineskip \topskip

%BibTex setup
\usepackage[semicolon,authoryear,square]{natbib}

\bibliographystyle{plainnat}
%%%%%%%%%%%%%%%%%%%%%%%%%%%%%%%%%%%%%%%%%%%%%%%%%%%%%%%%%%%%%%%%%%%%%%%%%%%%%%%%%%%%%%%%%%%%%%

% % Smart listing function. Documentation: http://texblog.org/tag/makenomenclature/
\usepackage{nomencl}
\makenomenclature 

% % % % % % % % % % % % % % % % Reference opsætning % % % % % % % % % % % % % % % % % % % % % % %

\usepackage{varioref} % for references, http://www.ctex.org/documents/packages/bibref/varioref.pdf
   
% External reference package, zref makes the references into hyperlinks.
\usepackage{url}
\usepackage{hyperref}
%\usepackage{zref-xr}
\usepackage{xr}

% % % % % % % % % % % % % % % % % % % % % % % % % % % % % % % % % % % % % % % % % % % % % % % % % % %

%\usepackage{showlabels}

% footruleskip is allready defined in memoir, so we have to undefine it, in order to include fancyhdr for our footer and header.
\let\footruleskip\undefined
\usepackage{fancyhdr}% For making fancy headers http://ctan.org/pkg/fancyhdr

\AtBeginDocument{\renewcommand\contentsname{Indholdsfortegnelse}}

% % % % % % % % % % % % % Inkluderinger af ønskede kapitler % % % % % % % % % % % % %

\begin{document}
\newgeometry{left=3cm,bottom=0.1cm}

\centerline{\Huge\huge\bfseries\color{ThemeColor}SMAP - Theme Project} 
\vspace{1em}
\centerline{\large\color{ThemeColor}Aarhus School of Engineering - Aarhus Universitet}


\centerline{\large\date{10. April 2015}}

%\color{Black}
%\vspace{8em} % Tilføj nogle tomme linjer
\begin{center}
    \begin{tabular}{ l l p{6cm} }
	 \textbf{Gruppe 5}& & \\
	& & \\
    Stud. nr.: 201270115 & Johnny Davidsen Poulsen & \\\hline
    & & \\
    Stud. nr.: 201270097  & Simon Østergaard Kristensen & \\\hline
    & & \\
    Stud. nr.: 201270108  & Michael Rahlff & \\\hline
    & & \\
    Stud. nr.: 201270763  & Sander Louring & \\\hline
     \end{tabular}
      	
\end{center}
\thispagestyle{empty} % No header and footer on this page
\restoregeometry
\setcounter{tocdepth}{2}
\tableofcontents

 

% Setup header and footer
\pagestyle{fancy}
\fancyhf{} % Clear all header an footer fields 
\renewcommand{\headrulewidth}{0pt} % Delete headruler
\fancyhead[R]{\includegraphics[height=40pt]{header}} % Insert picture in left side of header
%\fancyhead[R]{\rightmark} % insert chapter title in right side of header
\fancyfoot[LE,RO]{\thepage} % insert pagenumber on right side of even page, left side of odd page 
%\setcounter{page}{1}


\thispagestyle{fancy}
\chapter{Kravspecifikation}
\label{chp:description}

Med udgangspunkt i en agil udviklingsprocess opbygges kravene med user stories. De opbygges ud fra følgende skabelon: \\

\centerline{Som en <bruger> vil jeg <gøre noget> så <jeg kan opnå>}

\section{User stories}
\begin{enumerate}

\item Som en bruger, vil jeg kunne logge på med min Facebook, så at det bliver nemt for at oprette en bruger. 
\item Som en bruger, vil jeg introduceres til at et pengebeløb doneres baseret på min gåtur. 
\item Som en bruger, vil jeg kunne starte en gå tur.
\item Som en bruger, vil jeg kunne følge min gåtur på et kort. 
\item Som en bruger, vil jeg kunne følge antallet af kilometer jeg har gået. 
\item Som en bruger vil jeg kunne stoppe gåturen, når jeg er færdig. 
\item Som en bruger, vil jeg efter endt gåtur kunne se det donerede beløb baseret på min gåtur. 
\item Som en bruger, vil jeg efter endt gåtur kunne det samlede antal kilometer, som jeg er gået. 
\item Som bruger, vil jeg efter endt gåtur kunne dele min gåtur på Facebook, så jeg kan inspirere andre til at gøre det samme. 
\end{enumerate}
\thispagestyle{fancy}
\chapter{App design}
\label{chp:description}

\section{Sektion}

\thispagestyle{fancy}
\chapter{Software arkitektur}
\label{chp:description}

\section{Sektion}
\thispagestyle{fancy}
\chapter{Konklusion}
\label{chp:Konklusion}
Med udgangspunkt i kravspecifikationen, vil der i dette kapitel blive konkluderet på hvorvidt udviklingsprocessen og delimplementeringerne samlet opfylder kravspecifikationen. Konklusionen på opfyldelse af egen kravspecifikation vil blive suppleret med en refleksion over læringsprocessen og den tilegnede viden, i henhold til læringsmålene for faget og de formelle krav til smartphone applikationen.

Det er lykkedes i dette projekt at udarbejde en fuldt funktionel Android applikation som opfylder samtlige \emph{must} og \emph{should} krav, med en enkelt undtagelse. Eftersom at der i dette projekt udelukkende har været arbejdet med design og teknisk implementering, er det ikke lykkedes at opfylde kravet om at brugeren rent faktisk kan donere penge ved hjælp af sine gåture. For at dette krav fuldt ud skulle kunne indfries ville det både kræve at der efterfølgende blev indhentet sponsorer og at der blev udarbejdet en logistisk procedure for transaktioner af de donerede beløb.

I forhold til selve applikationen er vi nået i mål med samtlige user stories. Det er således muligt at logge ind i applikationen ved hjælp af sin Facebookkonto hvorefter man kan starte en gåtur og løbende se sin rute på et kort. Samtidig kan brugeren se hvor mange km., der er gået og hvor mange penge det svarer til i donationer. Når brugeren er færdig med at gå, kan brugeren afslutte turen, se sit resultat og dele det med andre brugere på Facebook. 

Der er desuden lavet et stykke forberedende arbejde i forhold til at gemme statistiske oplysninger centralt på en Google Cloud hostet service. Det er lykkedes at sende en HTTP anmodning fra vores applikation til en RESTful service, som vi selv har konfigureret. Det er ikke lykkedes at få denne service til at gemme de modtagende data i en SQL database, som vi har opsat ved Google.

\subsection{Refleksion}
Applikationen er på grund af projektperioden forløbstid udarbejdet gennem en enkel iteration, hvor vi har samarbejdet om en fælles kodebase i et git repository. Simple principper for continous integration er anvendt således at funktioner har kunne uddelegeres og delimplementeres. Det er således lykkedes at skabe en proces hvor integration har givet overraskende få problemer. Vi har arbejdet med Android frameworket i Android studio. Herudover er Google Services blevet anvendt til fremvisning af kort og til vores backend bestående af en RESTful service.


\thispagestyle{fancy}
\chapter{Arbejdsfordeling}
\label{chp:description}

Arbejdsfordeling beskrives med udgangspunkt i at man kan have påtaget sig en primær rolle og en sekundær rolle i arbejdet med en feature. Hele gruppen har deltaget ligeligt i arbejdet med rapporten og ligeledes design og arkitektur for applikationen.

\begin{table}[h]
\begin{tabular}{|>{\rr}p{3.2cm}|>{\rr}p{3.2cm}|>{\rr}p{3.2cm}|>{\rr}p{3.2cm}|p{3.2cm}|} 
\hline
\rowcolor{ThemeColor!80}  
\rowcolor{ThemeColor!80} \vspace{0.2cm} \textbf{Feature} \newline & \vspace{0.2cm}  \text{Johnny} & \vspace{0.2cm} \text{Rahlff} & \vspace{0.2cm}  \text{Sander} &  \vspace{0.2cm} \text{Simon} \\ 
\hline
 Google map tracking af rute  & P & S & S & P \vspace{1cm} \\
\hline 
Facebook integration & S & S &  P & P \vspace{1cm} \\
\hline
Backend og database & S & P & S & S \\
\hline
\end{tabular}
\caption{Arbejdsdeling}
\label{tab:MosCoW}
\end{table}
\FloatBarrier

%\bibliographystyle{plainnat}
%\bibliography{bibliografi}% Selects .bib file AND prints bibliography
\thispagestyle{fancy}

%\listoftodos % Make a list of todo's
\end{document}